\uuid{opLv}
\titre{Variables continues}
\auteur{worms}
\organisation{exo7}
\licence{CC BY-NC-SA 4.0}
\datecreate{2017-08-28}
\niveau{L2}
\module{Probabilité et statistique}
\chapitre{Probabilité continue}
\sousChapitre{Densité de probabilité}
\difficulte{}
\theme{Variables continues}
\exo7qcmid{worms.worms.046}

\contenu{
\begin{question}
\qclassification{261.01}
On suppose que $X$ désigne le montant du gain (en euros) que rapporte un employé à son entreprise en un mois, et que $X$ est de loi ${\cal N}(500,(150)^2)$. L'entreprise verse à l'employé une prime $Y$ égale à $0$ si ce gain est inférieur à 700 euros, et, si le gain $X$ excède 700 euros, à la moitié de l'excès en question.  Laquelle des affirmations suivantes est vraie ?
\begin{answers}
\bad{La loi de $Y$ est une loi continue et l'employé a environ $2,3\%$ de chances d'avoir une prime supérieure à 50 euros.}
\good{La loi de $Y$ n'est pas une loi continue et l'employé a environ $2,3\%$ de chances d'avoir une prime supérieure à 50 euros.}
\bad{La loi de $Y$ est une loi continue et l'employé a environ $17\%$ de chances d'avoir une prime supérieure à 50 euros.}
\bad{La loi de $Y$ n'est pas une loi continue et l'employé a environ $17\%$ de chances d'avoir une prime supérieure à 50 euros.}
\end{answers}
\begin{explanations}
On a $\{Y=0\} = \{X<700\}$, qui est de probabilité non nulle, ce qui contredit que $Y$ puisse être de loi continue (car on aurait alors $\mathbb{P}(Y=y)=0$ pour tout $y$ de $\Rr$). $Y$ n'est donc pas de loi continue, mais elle a une "composante" continue, on pourrait dire qu'elle continue conditionnellement au fait d'être $>0$. On a ensuite, pour $y >0$, $\{Y>y\} = \{Y>y \mbox{ , } X>700\} = \{ (X-700)/2 > y \mbox{ , } X>700\} = \{ X>700+2y \mbox{ , } X>700\} = \{X>700+2y\}$, donc $\{Y>50\} = \{X>800\} = \{ (X-500)/150 > (800-500)/150 \}$ qui est de probabilité $1-\Phi(2)$, qui vaut en effet environ $2.3\%$.  C'était une question plus difficile, mais très loin d'être inabordable ! Il faut juste écrire les événements tranquillement...
\end{explanations}
\end{question}
}
